\section{Three-Level Diagnostic Framework}

Endpoint metrics ask whether a latent representation contains task-relevant geometry; they do not test the distribution from which the planner actually selects actions. We therefore audit LeWM at three levels: fixed endpoint rankings, planner-produced candidate-pool rankings, and the selected rank-1 action.

\textbf{Endpoint ranking.} $\Rendpoint$ is the Spearman correlation between a terminal latent cost and the physical task cost on a fixed set of terminal-state/goal pairs. It measures whether the representation orders completed outcomes correctly, independent of CEM sampling and elite updates.

\textbf{Pool ranking.} $\Rpool$ is the same rank correlation, but computed inside the final 300-candidate pool produced by CEM for a start-goal pair. This is the ranking CEM actually uses at selection time. The CEM Compatibility Gap is
\[
\DeltaCEM = \Rendpoint - \Rpool.
\]
A large positive gap means endpoint geometry survives on a fixed benchmark but fails on the optimizer-local distribution.

\textbf{Selection metrics.} $M_\mathrm{rank1}$ records whether the lowest-cost candidate in the final pool succeeds when executed. We also report pool success mass and selection regret when available, because a pool can contain many successful actions even when rank ordering is weak.

The audit used fixed artifacts and pre-registered continuation gates to limit researcher freedom, but the process record is appendix material rather than the main argument; Appendix~\ref{app:preregistration} lists the decision trace, locked criteria, and git-anchored timeline.
